% Isolated proof note for the parent. Not independently audited or promoted.
% Parent owns the definitive all-odd scheme quotient/restriction theorem.
\section*{Proposed all-odd strengthening of the elliptic double-cover examples}

\textbf{Status: candidate proof note only.} This note records the complete
route for independent challenge after CP37. It makes no checkpoint or
canonical change and does not replace the already audited restricted
prime examples. The intended conclusion is exact geometric image
\(S_n\) and generic rank \(n-1\) for the same family at every odd
\(n=2m+1\ge3\).

Let \(E:y^2+y=x^3\) be supersingular over an algebraically closed
field of characteristic two. Put \(A=E\times E\), let \(q:A\to E\)
be the first projection, and use
\[
 V=\{O\}\times E,\quad H=E\times\{O\},\quad
 \Gamma=\{(x,y):x+y=0\},\quad
 L_0=\mathcal O_A(H+\Gamma),\quad L_m=L_0(mV).
\]
The already audited product-line lemma proves, for every \(m\ge1\),
that \(L_m\) is ample and globally generated, of type \((1,n)\)
with \(n=2m+1\), and that
\(\mathcal V_0=q_*L_0\) is a rank-two degree-one bundle.
In particular \(h^0(A,L_m)=n\). This input does not require a
prime degree or any generic Cartier-rank theorem.

\paragraph{The actual ruled quotient.}
On each fibre of \(q\), the restriction of \(L_0\) is the complete
degree-two line bundle \(\mathcal O_E(O+(-x))\). Fibrewise evaluation
therefore gives a morphism
\[
 \pi:A\longrightarrow S=\mathbf P_E(\mathcal V_0),\qquad
 \pi^*\mathcal O_S(1)=L_0,
\]
where the projective-bundle convention parametrizes one-dimensional
quotients. The map is proper and fibrewise finite of degree two, hence
finite; both surfaces are smooth, so it is flat. The degree-two elliptic
linear-system involution is
\[
                 \iota(x,y)=(x,-x-y).
\]
Thus \(\pi\) is the quotient by this involution. More explicitly,
it factors through the finite normal quotient \(A/\langle\iota\rangle\);
the resulting finite birational map to the normal surface \(S\) is
an isomorphism. This also supplies an actual descent of the line bundle,
not only numerical invariance under the involution.

\paragraph{The fixed elliptic curve and its quotient.}
The fixed scheme is the smooth elliptic curve
\[
                   T=\{(x,y):x=-[2]y\}\simeq E.
\]
Its tangent equation is \(dx=0\), since the differential of \([2]\)
vanishes. The differential of \(\iota-1\) has rank one at every
point of \(T\). Locally at a fixed point choose a parameter \(v\)
whose differential is nonzero in the vertical direction, and set
\(u=\iota(v)+v\). Then \(u,v\) are regular parameters,
\(\iota(u)=u\), \(\iota(v)=v+u\), and \(T=(u=0)\).
The completed invariant ring is exactly
\[
 k[[u,v]]^{\langle\iota\rangle}
                  =k[[u,w]],\qquad w=v^2+uv.
\]
Indeed \(k[[u,v]]\) is free of rank two over the displayed
subring, with basis \(1,v\); invariance of \(a+vb\) forces
\(ub=0\), hence \(b=0\). Consequently the reduced image
\(R=\pi(T)\) is smooth, and \(T\to R\) is a purely inseparable
degree-two map, locally \(w=v^2\). Thus \(R\) has genus one.
The original map \(q|_T=[-2]\) is purely inseparable of degree four,
because \(E\) is supersingular. Hence the induced \(R\to E\)
is purely inseparable of degree two as well. No geometric branch
points are identified by either map.

\paragraph{Restriction is a complete degree-\(n\) system.}
Set
\[
 \mathcal D_m=\mathcal O_S(1)\otimes\bar q^*\mathcal O_E(mO),
 \qquad N_m=\mathcal D_m|_R,
\]
where \(\bar q:S\to E\). All sections of \(L_m\) descend to
\(\mathcal D_m\): the pullback map on sections is an isomorphism,
as follows from the relative evaluation construction and
\(\bar q_*\mathcal D_m=\mathcal V_0(mO)=q_*L_m\).
Intersection on \(A\) gives
\[
 T\cdot V=4,\qquad T\cdot H=1,\qquad T\cdot\Gamma=1,
 \qquad L_m\cdot T=4m+2=2n.
\]
Since \(T\to R\) has degree two, \(\deg N_m=n\).
The restriction
\[
          H^0(S,\mathcal D_m)\longrightarrow H^0(R,N_m)
\]
is injective. A section in its kernel would pull back to a
nonzero section of \(L_m\) vanishing on \(T\), whereas
\((L_m-T)\cdot V=2-4=-2\) and \(V\) is nef, so
\(H^0(A,L_m-T)=0\). Both spaces have dimension \(n\), the
second by Riemann--Roch on the elliptic curve. Restriction is
therefore an actual linear isomorphism. Equivalently, restriction
to \(T\) consists of the pullbacks along \(T\to R\) of the
complete degree-\(n\) system, so locally these restrictions are
Frobenius squares. This conclusion uses an actual descended
bundle and quotient, rather than choosing arbitrary square roots
of individual sections.

\paragraph{The generic branch set and rank.}
For generic \(s\), the divisor \(D_s\) of \(\mathcal D_m\)
is a section of \(\bar q\). A section of the globally generated
rank-two bundle \(\mathcal V_0(mO)\) is generically nowhere zero:
the incidence of a zero at some point of the curve has codimension
at least one in its section space. Thus \(D_s\simeq E\).
Its pullback \(C_s\in|L_m|\) is smooth and connected on a
nonempty open by the already proved odd globally generated
smooth-member lemma. Independently, the complete degree-\(n\)
system on \(R\) has a nonempty open of reduced divisors, so
\(D_s\cap R\) consists of \(n\) distinct points generically.
These opens meet because the coefficient space is integral.

Off \(R\), the quotient \(\pi\) is etale. On \(R\), its
local equation is \(v^2+uv=w\). Thus the separable double cover
\(C_s\to D_s\simeq E\) has exactly the geometric branch
points \(D_s\cap R\). It is separable because \(D_s\) is not
contained in \(R\), so its quadratic equation has nonzero linear
coefficient in the generic field. The purely inseparable map
\(R\to E\) preserves the geometric branch permutation set,
with its Galois action. The canonical double-cover sheaf proof
already audited for the prime examples now gives
\[
 H^1_{\mathrm{et}}(C_{\bar\eta},\mathbf F_2)
          \simeq \mathbf F_2^{\mathcal B}/\mathbf F_2,
 \qquad |\mathcal B|=n.
\]
In particular the generic \(2\)-rank is \(n-1\), without invoking
the finite list of certified primes. Adjunction gives genus
\(n+1\), as before.

\paragraph{Exact symmetric branch monodromy.}
Under the restriction isomorphism, the branch family is the
universal reduced divisor of \(|N_m|\) on \(R\). Its ordered
pair incidence is a projective bundle over
\(R\times R\setminus\Delta\): two distinct evaluations are
independent since \(\deg N_m(-P-Q)=n-2>0\). This incidence
is irreducible, so the generic branch permutation group is
two-transitive.

Choose a divisor with one exact double point and every other
point simple. Such a divisor exists: for \(n\ge5\), the
residual degree \(n-2\) system is globally generated and its
general divisor is reduced away from the chosen point; for
\(n=3\), choose a point outside the finite set
\(N_m\simeq\mathcal O_R(3P)\). The latter set is finite since
\([3]\) is finite etale in characteristic two.

The local universal polynomial at the double point has, after
formal Weierstrass preparation, the form
\[
                       X^2+aX+b.
\]
The first two evaluation jets are independent, so \(a,b\)
are independent formal coefficient parameters. At the reduced
divisor \(a=0\), the element \(b\) is nonsquare in the
strict-henselian residue field; separably closed does not imply
perfect here. Adjoining a root gives a separable quadratic
extension of the generic field, with automorphism
\(X\mapsto X+a\) and a purely inseparable residue extension.
Its full inertia is therefore \(C_2\). All other roots are
etale and fixed after strict henselization. This supplies an
isolated transposition by fierce characteristic-two inertia;
it is not a tame-tangency argument. Two-transitivity and a
transposition imply the full group \(S_n\).

Finally its action on \(\mathbf F_2^n/\mathbf F_2\) is faithful
for odd \(n\ge3\): if a permutation acts trivially, each
\(e_i+e_{g(i)}\) is constant, but a nonzero such vector has
weight two, which is neither zero nor \(n\). Thus every index
is fixed. This would identify the actual relative geometric
Lang image with \(S_n\), and its nonzero orbits with the
complement-paired subset weights \(1,\ldots,(n-1)/2\).

\paragraph{Scope for the independent audit.}
The essential new steps are the global identification of the
relative degree-two quotient, the completed invariant-ring
calculation along its fixed curve, the actual complete-system
restriction, and the local fierce inertia. They must all be
checked before this note is promoted. No numerical test,
characteristic-zero scroll theorem, or general fixed-embedding
full-monodromy conclusion is asserted. Parent ownership of the
definitive theorem prevents duplicate shared integration.
